\documentclass[11pt]{article}
\usepackage[letterpaper,margin=0.68in]{geometry}
\usepackage{amsmath,amssymb}
\usepackage{enumitem}
\usepackage{xcolor}
\usepackage{fancyhdr}
\usepackage{microtype}
\usepackage{tabularx,array,booktabs}
\usepackage{tikz}
\usepackage[most]{tcolorbox}

\definecolor{olive}{HTML}{4D5430}
\definecolor{gold}{HTML}{9A7B22}
\definecolor{pale}{HTML}{F5F1DC}
\definecolor{softgreen}{HTML}{EDF0E2}
\definecolor{softgray}{HTML}{F5F5F2}
\definecolor{ink}{HTML}{292D20}

\pagestyle{fancy}
\fancyhf{}
\lhead{\textcolor{olive}{MATH\& 107: Math in Society}}
\rhead{\textcolor{gold}{Fall 2026}}
\cfoot{\thepage}
\setlength{\headheight}{14pt}
\setlength{\parindent}{0pt}
\setlength{\parskip}{5pt}
\setlist[enumerate]{leftmargin=*,itemsep=7pt,topsep=5pt}
\setlist[itemize]{leftmargin=1.35em,itemsep=3pt,topsep=3pt}

\newtcolorbox{keybox}[1]{colback=pale,colframe=olive,boxrule=0.8pt,arc=2mm,title=\textbf{#1},fonttitle=\color{olive},left=2mm,right=2mm,top=1.5mm,bottom=1.5mm}
\newtcolorbox{examplebox}[1]{colback=softgreen,colframe=gold,boxrule=0.8pt,arc=2mm,title=\textbf{#1},fonttitle=\color{ink},left=2mm,right=2mm,top=1.5mm,bottom=1.5mm}
\newtcolorbox{refbox}[1]{colback=softgray,colframe=gold,boxrule=0.7pt,arc=2mm,title=\textbf{#1},fonttitle=\color{ink},left=2mm,right=2mm,top=1.5mm,bottom=1.5mm}

\newcommand{\summaryhead}[2]{%
  \clearpage
  {\Large\bfseries\textcolor{olive}{Module #1: #2 — Summary}}\par
  \vspace{2pt}\textcolor{gold}{\rule{\linewidth}{1.2pt}}\par\smallskip
}
\newcommand{\prequizhead}[2]{%
  \clearpage
  {\Large\bfseries\textcolor{olive}{Module #1: #2 — Prequiz}}\par
  \textit{Open-note. Show your mathematical work. A calculation and a clearly stated final answer are enough unless a sentence is specifically requested.}\par
  Name: \hspace{2.6in} Date: \hspace{1.1in}\par
  \vspace{2pt}\textcolor{gold}{\rule{\linewidth}{1.2pt}}\par\smallskip
}
\newcommand{\quizhead}[2]{%
  \clearpage
  {\Large\bfseries\textcolor{olive}{Module #1: #2 — Matching Quiz}}\par
  \textit{These problems use the same skills as the prequiz, with new values and contexts. Show enough work to make your reasoning visible.}\par
  Name: \hspace{2.6in} Date: \hspace{1.1in}\par
  \vspace{2pt}\textcolor{gold}{\rule{\linewidth}{1.2pt}}\par\smallskip
}
\newcommand{\worksmall}{\par\vspace{0.38in}}
\newcommand{\workmed}{\par\vspace{0.62in}}
\newcommand{\worklarge}{\par\vspace{0.92in}}
\newcommand{\workxl}{\par\vspace{1.22in}}
\newcommand{\choice}[1]{\(\square\)~#1\qquad}

\begin{document}
\summaryhead{6}{Interest and Loans}

\textbf{What this module is about.} Interest connects arithmetic growth to real financial decisions. Simple interest adds the same dollar amount each period because it is always computed from the original principal. Compound interest multiplies the current balance by the same growth factor, so the dollar increase itself grows over time. Loans use the same percentage-growth idea, but a payment is subtracted after interest is added.

\begin{keybox}{Simple and compound interest}
If \(P\) is principal, \(r\) is the annual rate as a decimal, and \(t\) is time in years, simple interest is
\[
A=P(1+rt).
\]
If interest is compounded \(n\) times per year,
\[
A=P\left(1+\frac rn\right)^{nt}.
\]
Simple interest has a constant dollar increase. Compound interest has a constant multiplicative factor per compounding period.
\end{keybox}

\begin{examplebox}{Example 1: Simple versus compound growth}
Invest \(\$1500\) at \(6\%\) for \(10\) years.

Simple interest:
\[
A=1500(1+0.06\cdot10)=1500(1.6)=\$2400.
\]
Annual compounding:
\[
A=1500(1.06)^{10}\approx\$2686.27.
\]
The compound balance is larger because later interest is earned on earlier interest.
\end{examplebox}

\begin{keybox}{Rule of 72}
For percentage growth near ordinary financial rates, the doubling time can be estimated by
\[
t\approx\frac{72}{r_{\%}},
\]
where \(r_{\%}\) is the annual percentage rate written as a percent number. At \(6\%\), the estimated doubling time is \(72/6=12\) years. To estimate a quadrupling time, remember that quadrupling is two doublings.
\end{keybox}

\begin{keybox}{Loans: interest first, then payment}
For a loan balance \(B\), annual rate \(r\), and monthly payment \(M\), the monthly rate is
\[
i=\frac r{12}.
\]
During one month:
\[
\text{interest}=iB,
\qquad
\text{principal reduction}=M-iB,
\]
\[
B_{\text{new}}=B(1+i)-M.
\]
If the payment is not larger than the interest charge, the balance will not decrease.
\end{keybox}

\begin{examplebox}{Example 2: One month of a loan}
A \(\$3000\) loan has \(9\%\) APR compounded monthly and a \(\$150\) payment.
\[
i=\frac{0.09}{12}=0.0075.
\]
First-month interest:
\[
0.0075(3000)=\$22.50.
\]
Of the \(\$150\) payment, \(150-22.50=\$127.50\) reduces principal. The new balance is
\[
3000+22.50-150=\$2872.50.
\]
A larger payment on the same starting balance does not change the first month's interest; it changes how much principal is removed and therefore changes future interest.
\end{examplebox}

\textbf{By the end of this module, you should be able to:} calculate simple and compound balances; identify the number of compounding periods and the periodic rate; use the Rule of 72; and perform a one-period loan update showing interest, principal reduction, and new balance.

\prequizhead{6}{Interest and Loans}
\begin{refbox}{Useful formulas}
\[
A=P(1+rt),
\qquad
A=P\left(1+\frac rn\right)^{nt},
\qquad
t_{\text{double}}\approx\frac{72}{r_{\%}},
\]
\[
B_{\text{new}}=B\left(1+\frac r{12}\right)-M.
\]
\end{refbox}

\begin{enumerate}
\item A student deposits \(\$1500\) in an account that earns \(4\%\) simple interest per year for \(3\) years.
\begin{enumerate}[label=(\alph*),itemsep=4pt]
\item Find the total amount in the account after \(3\) years.
\item How many dollars of interest are added each year under this simple-interest model?
\end{enumerate}
\worklarge

\item A savings account starts with \(\$1000\) and earns \(5\%\) compounded annually for \(4\) years.
\begin{enumerate}[label=(\alph*),itemsep=4pt]
\item Write the compound-interest expression.
\item Find the balance to the nearest cent.
\item Explain in one short phrase why the dollar amount of interest is not the same every year.
\end{enumerate}
\worklarge

\item Use the Rule of 72 for an investment growing at approximately \(6\%\) per year.
\begin{enumerate}[label=(\alph*),itemsep=4pt]
\item Estimate the doubling time.
\item Estimate the time required to grow to roughly four times the original amount.
\end{enumerate}
\workmed

\item A \(\$3000\) loan charges \(9\%\) APR compounded monthly. The borrower pays \(\$150\) at the end of each month.
\begin{enumerate}[label=(\alph*),itemsep=4pt]
\item Find the monthly rate as a decimal.
\item Find the interest charged during the first month.
\item How much of the first payment reduces principal?
\item Find the balance immediately after the first payment.
\end{enumerate}
\workxl

\item Two loans have the same starting balance and the same interest rate. Loan A has a monthly payment of \(\$150\), while Loan B has a monthly payment of \(\$250\). Which loan has the larger first-month interest charge? Which loan will generally be paid off sooner? Explain each answer using the loan update process rather than guessing from the payment size alone.
\worklarge
\end{enumerate}

\quizhead{6}{Interest and Loans}
\begin{enumerate}
\item A \(\$2000\) account earns \(3\%\) simple interest for \(5\) years. Find the final balance and the dollar amount added each year.
\workmed

\item A \(\$500\) account earns \(8\%\) compounded annually for \(3\) years. Find the final balance to the nearest cent.
\workmed

\item Use the Rule of 72 to estimate the doubling time at \(8\%\) annual growth. About how long would two doublings take?
\workmed

\item A \(\$6000\) loan charges \(12\%\) APR compounded monthly and has a \(\$200\) payment. Find the monthly rate, first-month interest, amount of principal reduction, and new balance after the payment.
\worklarge

\item Loan A and Loan B both begin at \(\$5000\) with the same APR. Loan A requires \(\$150\) each month and Loan B requires \(\$250\). Which has the same first-month interest? Which payment removes more principal in month 1?
\workmed
\end{enumerate}

% =====================================================================
% MODULE 7
% =====================================================================
\end{document}
